\documentclass[../main.tex]{subfiles}
\begin{document}

\subsection{Равномерное движение по окружности. Угловая скорость}
%==============================================================================

Окружность~--- простейшая криволинейная траектория, на которой все
свойства кривого движения проявляются особенно наглядно.
Начнём с равномерного движения и постепенно перейдём к более общим случаям.

\noindent\textbf{Угловая скорость и период.}
Пусть материальная точка движется по окружности радиуса $R$ с постоянной
по величине скоростью $v$. За малое время $\Delta t$ она проходит дугу
$\Delta s = v\,\Delta t$ и поворачивается на угол $\Delta\varphi = \frac{\Delta s}{R}$.
\textit{Угловой скоростью} называется скорость изменения этого угла:
\[
\omega = \lim_{\Delta t \to 0}\frac{\Delta\varphi}{\Delta t} = \frac{v}{R}.
\]
Время одного полного оборота ($\Delta\varphi = 2\pi$) называется \textit{периодом}:
\[
T = \frac{2\pi}{\omega} = \frac{2\pi R}{v}.
\]

\noindent\textbf{Уравнения движения.}
Поместим центр окружности в начало координат; в момент $t = 0$ точка
находится в точке $(R, 0)$.
Через время $t$ она повернулась на угол $\omega t$, и её координаты:
\[
x(t) = R\cos(\omega t), \qquad y(t) = R\sin(\omega t).
\]
Компоненты скорости находятся как пределы отношений приращений координат
к $\Delta t$. Для равномерного движения они принимают вид:
\[
v_x = -R\omega\sin(\omega t), \qquad
v_y = \phantom{-}R\omega\cos(\omega t).
\]
Можно убедиться, что $v_x^2 + v_y^2 = R^2\omega^2 = v^2$~--- скорость
по величине постоянна. Проверим также перпендикулярность скорости и радиус-вектора:
$\vec{v}\cdot\vec{r} = (-R\omega\sin\omega t)(R\cos\omega t)
+ (R\omega\cos\omega t)(R\sin\omega t) = 0$.

\noindent\textbf{Центростремительное ускорение.}
Скорость постоянна по величине, поэтому меняется только её направление.
За малое время $\Delta t$ вектор скорости повернулся на тот же угол
$\Delta\varphi = \frac{v\,\Delta t}{R}$, что и радиус-вектор
(скорость всегда перпендикулярна ему).
Для малого угла поворота модуль изменения скорости:
\[
|\Delta\vec{v}| \approx v \cdot \Delta\varphi = v \cdot \frac{v\,\Delta t}{R}
= \frac{v^2}{R}\Delta t.
\]
Поделим на $\Delta t$ и получим центростремительное ускорение:
\[
a = \lim_{\Delta t \to 0}\frac{|\Delta\vec{v}|}{\Delta t} = \frac{v^2}{R} = \omega^2 R.
\]
Направление $\Delta\vec{v}$~--- перпендикулярно скорости, то есть
к центру окружности. Именно это ускорение непрерывно поворачивает
вектор скорости, не давая точке улететь по прямой.

\noindent\textbf{Угловая скорость радиус-вектора.}
Полезно обобщить понятие угловой скорости на произвольную кривую.
Пусть точка движется по любой траектории. В каждый момент только
та компонента скорости, которая перпендикулярна радиус-вектору $\vec{r}$,
вызывает вращение этого вектора. Радиальная компонента лишь изменяет
длину $r = |\vec{r}|$. Поэтому угловая скорость радиус-вектора:
\[
\omega_r = \frac{v_\perp}{r},
\]
где $v_\perp$~--- проекция скорости, перпендикулярная $\vec{r}$.
Для равномерного движения по окружности $v_\perp = v$ и $r = R$,
откуда $\omega_r = \frac{v}{R} = \omega$.

\begin{example}
Материальная точка движется по окружности с постоянной скоростью
$v = 2$ м/с. За время $\tau = 2$ с вектор скорости повернулся на угол
$\alpha = 45^\circ$. Найдите период обращения и центростремительное
ускорение точки.
При равномерном движении вектор скорости поворачивается с той же угловой
скоростью $\omega$, что и сама точка. Из условия:
$\omega = \frac{\alpha}{\tau} = \frac{\pi/4}{2} = \frac{\pi}{8}$~рад/с.
Период: $T = \frac{2\pi}{\omega} = 16$~с.
Радиус: $R = \frac{v}{\omega} = \frac{16}{\pi}$~м.
Центростремительное ускорение:
\[
a = v\omega = 2 \cdot \frac{\pi}{8} = \frac{\pi}{4} \approx 0{,}785~\text{м/с}^2.
\]
\end{example}

\begin{example}
Тело движется по окружности радиуса $R$ с угловой скоростью $\omega$.
Центр окружности совпадает с началом координат; в момент $t = 0$
тело находится на оси $OX$ с положительной координатой.
Найдите угловую скорость вращения отрезка, соединяющего тело
с точкой $P = (-R,\, 0)$.
Положение тела: $A = (R\cos\omega t,\, R\sin\omega t)$.
Вектор $\overrightarrow{PA} = (R\cos\omega t + R,\, R\sin\omega t)$.
Воспользуемся формулами двойного угла
$1 + \cos\theta = 2\cos^2(\theta/2)$ и $\sin\theta = 2\sin(\theta/2)\cos(\theta/2)$:
\[
\overrightarrow{PA} = 2R\cos\frac{\omega t}{2}
\cdot \left(\cos\frac{\omega t}{2},\; \sin\frac{\omega t}{2}\right).
\]
Это вектор, направленный под углом $\frac{\omega t}{2}$ к оси $OX$.
Угловая скорость его вращения:
\[
\Omega = \frac{\Delta(\omega t/2)}{\Delta t} = \frac{\omega}{2}.
\]
Отрезок $PA$ вращается вдвое медленнее, чем само тело.
Этот результат~--- аналог теоремы о вписанном угле: угол, вписанный
в окружность, вдвое меньше центрального, опирающегося на ту же дугу.
\end{example}

%==============================================================================
\subsection{Нормальное и касательное ускорения}
%==============================================================================

При равномерном движении по окружности ускорение было целиком центростремительным.
Теперь рассмотрим неравномерное движение, при котором
скорость меняется и по величине, и по направлению.

\noindent\textbf{Векторный треугольник скоростей.}
В момент $t$ вектор скорости равен $\vec{v}$, а в момент $t + \Delta t$~---
$\vec{v}'$. Вектор приращения $\Delta\vec{v} = \vec{v}' - \vec{v}$
замыкает треугольник, образованный $\vec{v}$ и $\vec{v}'$.
Разложим $\Delta\vec{v}$ на две составляющие: \textit{вдоль} $\vec{v}$
и \textit{перпендикулярно} ей.

\textit{Составляющая вдоль $\vec{v}$} равна изменению модуля скорости:
$\Delta v_\tau \approx |\vec{v}'| - |\vec{v}|$.
В пределе $\Delta t \to 0$ это даёт \textit{касательное ускорение}:
\[
a_\tau = \lim_{\Delta t \to 0}\frac{\Delta v_\tau}{\Delta t}.
\]

\textit{Составляющая перпендикулярная $\vec{v}$} связана
с поворотом скорости на угол $\Delta\varphi$.
При малом $\Delta\varphi$ её величина $\Delta v_n \approx v\,\Delta\varphi$.
В пределе:
\[
a_n = \lim_{\Delta t \to 0}\frac{v\,\Delta\varphi}{\Delta t} = v\,\omega_v,
\]
где $\omega_v = \lim_{\Delta t \to 0}\frac{\Delta\varphi}{\Delta t}$~--- угловая скорость вращения вектора скорости.
Для точки на окружности радиуса $R$: $\Delta\varphi = \frac{\Delta s}{R} = \frac{v\,\Delta t}{R}$,
поэтому $\omega_v = \frac{v}{R}$ и $a_n = \frac{v^2}{R}$.

\begin{definition}
Ускорение точки, движущейся по криволинейной траектории,
раскладывается на два взаимно перпендикулярных слагаемых:
\[
a_\tau = \lim_{\Delta t \to 0}\frac{\Delta v}{\Delta t}
\qquad\text{--- \textit{касательное ускорение},}
\]
\[
a_n = \frac{v^2}{\rho}
\qquad\text{--- \textit{нормальное (центростремительное) ускорение}.}
\]
Касательное ускорение отвечает за изменение быстроты движения;
нормальное~--- за изменение направления скорости.
Полное ускорение и угол $\theta$ между ним и скоростью:
\[
|\vec{a}|^2 = a_\tau^2 + a_n^2,
\qquad
\tan\theta = \frac{a_n}{a_\tau}.
\]
\end{definition}

При равномерном движении ($a_\tau = 0$) ускорение целиком направлено к центру;
при прямолинейном ($a_n = 0$)~--- вдоль скорости.

\begin{example}
Частица начинает двигаться по окружности радиуса $R$ из состояния
покоя с постоянным касательным ускорением $a_\tau$.
Найдите угол $\theta$ между вектором ускорения и вектором скорости
после первого оборота.
После оборота точка прошла путь $s = 2\pi R$.
Скорость из кинематики (при постоянном $a_\tau$):
$v^2 = 2a_\tau s = 4\pi R a_\tau$.
Нормальное ускорение: $a_n = \frac{v^2}{R} = 4\pi a_\tau$.
\[
\tan\theta = \frac{a_n}{a_\tau} = 4\pi
\quad\Rightarrow\quad
\theta = \arctan(4\pi) \approx 85{,}5^\circ.
\]
Полное ускорение почти перпендикулярно скорости: нормальная составляющая
в $4\pi \approx 12{,}6$ раза больше касательной.
\end{example}

\begin{example}
Точка движется по окружности радиуса $R$ из состояния покоя,
при этом модуль полного ускорения постоянен и равен $a$.
Какую часть окружности пройдёт точка к моменту, когда она
достигнет максимальной скорости?
Максимальная скорость достигается, когда касательное ускорение
обращается в нуль: все ускорение становится нормальным ($a_n = a$).
Заметим, что угол $\varphi$ между вектором ускорения и скоростью
начинается с нуля (в начальный момент $v = 0$, ускорение целиком касательное)
и заканчивается значением $\frac{\pi}{2}$ (в момент максимальной скорости).
\textit{Ключевое утверждение:} угол $\varphi$ изменяется вдвое быстрее,
чем угол поворота точки на окружности. Строгое доказательство:
пусть $\omega = \frac{v}{R}$~--- угловая скорость точки. Из соотношений
$\cos\varphi = \frac{a_\tau}{a}$ и $a\sin\varphi = a_n = \frac{v^2}{R}$ рассмотрим малые приращения.
Приращение левой части второго равенства равно $a\cos\varphi\,\Delta\varphi$,
приращение правой~--- $\frac{2v}{R}\Delta v$.
Приравнивая их и деля на $\Delta t$, получаем
\[
a\cos\varphi\,\frac{\Delta\varphi}{\Delta t} = \frac{2v}{R}\frac{\Delta v}{\Delta t}.
\]
Учитывая, что $\frac{\Delta v}{\Delta t} = a_\tau = a\cos\varphi$, сокращаем и получаем
\[
\frac{\Delta\varphi}{\Delta t} = \frac{2v}{R} = 2\omega.
\]
Итак, угловая скорость изменения угла наклона ускорения в два раза больше
угловой скорости вращения точки. Изменение угла поворота $\theta$ при переходе
$\varphi$ от $0$ до $\frac{\pi}{2}$:
\[
\Delta\theta = \frac{\Delta\varphi}{2} = \frac{\pi/2}{2} = \frac{\pi}{4}.
\]
Доля окружности: $\frac{\Delta\theta}{2\pi} = \frac{1}{8}$.
\end{example}

%==============================================================================
\subsection{Радиус кривизны}
%==============================================================================

До сих пор нормальное ускорение вычислялось по формуле $a_n = \frac{v^2}{\rho}$,
где $\rho$~--- радиус соприкасающейся окружности. Но если траектория~--- не окружность?
В каждой точке произвольной кривой изогнутость может быть различной.
Нам нужен инструмент, позволяющий измерить эту изогнутость в отдельной точке.

\noindent\textbf{Соприкасающаяся окружность.}
Через любые три точки, не лежащие на одной прямой, проходит
единственная окружность (описанная окружность треугольника).
Зафиксируем на кривой точку $P$ и выберем ещё две точки $Q$ и $S$
по обе стороны от неё. Проведём описанную окружность через $P$, $Q$, $S$.
При стягивании $Q$ и $S$ к точке $P$ вдоль кривой эта окружность
стремится к пределу~--- \textit{соприкасающейся окружности}.
Она «наилучшим образом» приближает кривую в точке $P$:
совпадает с ней до второго порядка (общая точка, общее направление
и одинаковая кривизна). Геометрически это означает, что вблизи $P$
кривая и окружность имеют одинаковый наклон и одинаковую скорость
поворота касательной.

\begin{definition}
\textit{Радиусом кривизны} $\rho$ в точке $P$ называется радиус
соприкасающейся окружности. Чем сильнее кривая изогнута в данной точке,
тем меньше $\rho$. Для прямой линии $\rho = \infty$.
\end{definition}

\noindent\textbf{Связь с ускорением.}
Для точки, движущейся по произвольной кривой, нормальное ускорение
в каждый момент времени~--- это центростремительное ускорение на
соприкасающейся окружности радиуса $\rho$:
\[
a_n = \frac{v^2}{\rho}.
\]
Отсюда получаем рабочую формулу: $\rho = \frac{v^2}{a_n}$.
Зная скорость точки и нормальное ускорение, можно найти радиус кривизны
траектории в данный момент~--- даже не зная уравнения самой траектории.

\noindent\textbf{Угловая скорость вращения скорости.}
Из формулы $a_n = v\,\omega_v$ и $a_n = \frac{v^2}{\rho}$ следует:
\[
\omega_v = \frac{v}{\rho}.
\]
Чем меньше радиус кривизны (резче поворот), тем быстрее поворачивается
вектор скорости.

\noindent\textbf{Аналитическое вычисление радиуса кривизны.}
До сих пор $\rho$ находился из кинематических величин $v$ и $a_n$.
Однако $\rho$~--- это прежде всего \textit{геометрическое} свойство самой кривой,
не зависящее от закона движения точки по ней. Выведем формулу,
позволяющую вычислять $\rho$ непосредственно из уравнения траектории $y = f(x)$.

Для кривой $y = f(x)$ касательная в каждой точке образует
с осью $OX$ угол $\varphi$, задаваемый соотношением:
\[
    \tan\varphi = y' \equiv \frac{\dd y}{\dd x}.
\]
Малый элемент длины дуги $ds$ находится по теореме Пифагора:
\[
    ds = \sqrt{dx^2 + dy^2} = \sqrt{1 + (y')^2}\;dx,
    \qquad\text{откуда}\quad
    \frac{ds}{dx} = \sqrt{1 + (y')^2}.
\]
Рассматривая прямоугольный треугольник с катетами $dx$, $dy$ и гипотенузой $ds$,
получаем геометрическое соотношение для косинуса угла наклона:
\[
    \cos\varphi = \frac{dx}{ds} = \frac{1}{\sqrt{1 + (y')^2}}.
\]
Степень изогнутости кривой в точке естественно измерять тем, насколько быстро
поворачивается касательная при продвижении вдоль кривой. Это приводит к
определению \textit{кривизны}~--- производной угла касательной по длине дуги:
\[
    \kappa = \left|\frac{\dd\varphi}{\dd s}\right|.
\]
Радиус кривизны и кривизна обратно пропорциональны: $\rho = \frac{1}{\kappa}$.
Для прямой линии угол $\varphi$ постоянен, поэтому $\kappa = 0$ и $\rho = \infty$.
Для окружности радиуса $R$ касательная поворачивается равномерно:
при прохождении дуги $ds$ угол меняется на $d\varphi = ds/R$, откуда
$\kappa = \frac{1}{R}$ и $\rho = R$, что согласуется с геометрической интуицией.

Выразим $\frac{\dd\varphi}{\dd s}$ через производные $y'$ и $y''$.
Дифференцируем равенство $\tan\varphi = y'$ по $x$, используя правило
производной сложной функции:
\[
    \frac{1}{\cos^2\varphi}\cdot\frac{\dd\varphi}{\dd x} = y''
    \quad\Longrightarrow\quad
    \frac{\dd\varphi}{\dd x} = y''\cos^2\varphi.
\]
Чтобы перейти от производной по $x$ к производной по длине дуги,
применим правило дифференцирования сложной функции $\frac{\dd\varphi}{\dd s} = \frac{\dd\varphi/\dd x}{\dd s/\dd x}$:
\[
    \frac{\dd\varphi}{\dd s}
    = \frac{y''\cos^2\varphi}{\sqrt{1 + (y')^2}}
    = \frac{y''\cos^2\varphi}{1/\cos\varphi}
    = y''\cos^3\varphi.
\]
Подставляем выражение $\cos\varphi = \frac{1}{\sqrt{1+(y')^2}}$:
\[
    \frac{\dd\varphi}{\dd s} = \frac{y''}{\bigl(1 + (y')^2\bigr)^{3/2}}.
\]
Взяв модуль и используя связь $\rho = \left|\frac{\dd s}{\dd\varphi}\right|$, получаем окончательный результат:
\[
    \boxed{\rho = \frac{\bigl(1 + (y')^2\bigr)^{3/2}}{|y''|}}.
\]

\noindent\textbf{Связь с кинематической формулой.}
Убедимся, что геометрический вывод согласуется с формулой $\rho = \frac{v^2}{a_n}$.
Пусть точка движется по кривой $y = f(x)$ так, что горизонтальная
компонента скорости постоянна: $\dot x = u$. Тогда $\dot y = y'u$,
$\ddot y = y''u^2$, $\ddot x = 0$, а модуль скорости $v = u\sqrt{1+(y')^2}$.
В двумерном случае модуль векторного произведения скорости и ускорения равен
$|\vec{v}\times\vec{a}| = |\dot x\,\ddot y - \dot y\,\ddot x|$. Поскольку
нормальное ускорение связано с этим произведением соотношением
$a_n = \frac{|\vec{v}\times\vec{a}|}{v}$, получаем:
\[
    a_n = \frac{|u \cdot y'' u^2 - y'u\cdot 0|}{u\sqrt{1+(y')^2}}
        = \frac{|y''|u^2}{\sqrt{1+(y')^2}}.
\]
Вычисляем радиус кривизны кинематически:
\[
    \rho = \frac{v^2}{a_n} = \frac{u^2\bigl(1+(y')^2\bigr)}{|y''|u^2/\sqrt{1+(y')^2}}
         = \frac{\bigl(1+(y')^2\bigr)^{3/2}}{|y''|}.
\]
Геометрический и кинематический подходы приводят к идентичному выражению.

\noindent\textbf{Параметрическая запись.}
Если кривая задана параметрически $x = x(t)$, $y = y(t)$ (параметр $t$
не обязательно совпадает со временем), аналогичные выкладки дают:
\[
    \rho = \frac{\bigl(\dot x^2 + \dot y^2\bigr)^{3/2}}{|\dot x\,\ddot y - \dot y\,\ddot x|}.
\]
Здесь числитель равен $v^3$, а знаменатель~--- $v\,a_n$, поэтому частное
в точности воспроизводит фундаментальное соотношение $\rho = \frac{v^2}{a_n}$.

\begin{example}
Проверим формулу для окружности $x^2 + y^2 = R^2$.
Рассмотрим верхнюю полуокружность $y = \sqrt{R^2 - x^2}$. Находим производные:
\[
    y' = -\frac{x}{\sqrt{R^2 - x^2}},
    \qquad
    y'' = -\frac{R^2}{(R^2 - x^2)^{3/2}}.
\]
Вычисляем выражение в числителе:
\[
    1 + (y')^2 = 1 + \frac{x^2}{R^2 - x^2} = \frac{R^2}{R^2 - x^2}.
\]
Подставляем в общую формулу (знак второй производной учитывается модулем):
\[
    \rho = \frac{\bigl(R^2/(R^2-x^2)\bigr)^{3/2}}{R^2/(R^2-x^2)^{3/2}}
         = \frac{R^3/(R^2-x^2)^{3/2}}{R^2/(R^2-x^2)^{3/2}} = R.
\]
Радиус кривизны окружности постоянен и равен её геометрическому радиусу,
что и требовалось проверить.
\end{example}

\begin{example}
Найдём радиус кривизны траектории тела, брошенного горизонтально со
скоростью $v_0$. Уравнение траектории имеет вид параболы:
$y = \frac{g}{2v_0^2}\,x^2$.

Вычисляем производные:
\[
    y' = \frac{gx}{v_0^2},
    \qquad
    y'' = \frac{g}{v_0^2}.
\]
Подставляем в формулу:
\[
    \rho(x) = \frac{\Bigl(1 + \frac{g^2x^2}{v_0^4}\Bigr)^{3/2}}{\frac{g}{v_0^2}}
            = \frac{v_0^2}{g}\left(1 + \frac{g^2x^2}{v_0^4}\right)^{3/2}.
\]
В начальный момент ($x = 0$) получаем $\rho_{\min} = \frac{v_0^2}{g}$ ---
наименьший радиус кривизны, соответствующий максимальной изогнутости
траектории в вершине параболы. По мере удаления от точки броска
($x \to \infty$) кривая распрямляется и $\rho(x) \to \infty$.
Этот результат полностью совпадает с ответом, полученным ранее
кинематическим методом через нормальное ускорение.
\end{example}

%==============================================================================
\subsection{Задачи преследования}
%==============================================================================

В задачах преследования одно тело («преследователь») всегда движется
в направлении другого («цели»). Скорость при этом постоянна по величине,
но непрерывно меняет направление. Ключевым инструментом служит
производная расстояния между преследователем и целью.

\noindent\textbf{Скорость сближения.}
Пусть $A$~--- преследователь, $B$~--- цель. Обозначим $r = |AB|$.
Изменение расстояния за малое время $\Delta t$~--- это проекция вектора
относительной скорости $\vec{v}_B - \vec{v}_A$ на прямую $AB$:
\[
\Delta r = \vec{e}_{AB}\cdot(\vec{v}_B - \vec{v}_A)\,\Delta t,
\]
где $\vec{e}_{AB}$~--- единичный вектор от $A$ к $B$.
Если $\Delta r < 0$, расстояние убывает (тела сближаются).
Раскроем явно: пусть $\beta_A$~--- угол между $\vec{v}_A$ и направлением $AB$,
а $\beta_B$~--- угол между $\vec{v}_B$ и тем же направлением. Тогда:
\[
\frac{\Delta r}{\Delta t} = v_B\cos\beta_B - v_A\cos\beta_A.
\]

\begin{example}
\textbf{Черепахи в вершинах равностороннего треугольника.}
Три черепахи находятся в вершинах правильного треугольника со стороной $a$.
Каждая движется к своей соседке по часовой стрелке с постоянной скоростью $v$.
Найдите время встречи, путь каждой черепахи и место встречи.

Начальная конфигурация и законы движения инвариантны относительно поворота на $120^\circ$
вокруг центра треугольника. Запишем уравнение движения для $k$-й черепахи:
\[
\frac{d\vec{r}_k}{dt} = v\,\frac{\vec{r}_{k+1} - \vec{r}_k}{|\vec{r}_{k+1} - \vec{r}_k|}.
\]
Поскольку правая часть преобразуется согласованно с циклической перестановкой индексов,
симметрия положений сохраняется во времени: черепахи всегда образуют правильный треугольник,
который равномерно уменьшается в размерах и вращается. Центр треугольника остаётся
неподвижным, и именно в нём произойдёт встреча.

Для нахождения времени встречи вычислим скорость изменения расстояния $r(t)$ между
двумя соседними черепахами $A$ и $B$. Используем общую формулу для производной расстояния:
\[
\frac{\Delta r}{\Delta t} = v_B\cos\beta_B - v_A\cos\beta_A,
\]
где $\beta_A$ --- угол между $\vec{v}_A$ и направлением $AB$, $\beta_B$ --- угол между
$\vec{v}_B$ и тем же направлением. По условию $\vec{v}_A$ направлен строго к $B$,
поэтому $\beta_A = 0$, $\cos\beta_A = 1$. Черепаха $B$ движется к $C$; угол
между направлением $BC$ и прямой $AB$ равен внешнему углу правильного треугольника,
то есть $120^\circ$, следовательно $\cos\beta_B = -\frac{1}{2}$.
Подставляем в формулу:
\[
\frac{\Delta r}{\Delta t} = v\left(-\frac{1}{2}\right) - v(1) = -\frac{3v}{2}.
\]
Скорость сближения постоянна: $v_{\text{сбл}} = -\frac{\Delta r}{\Delta t} = \frac{3v}{2}$.
Время встречи:
\[
t = \frac{a}{v_{\text{сбл}}} = \frac{2a}{3v}.
\]
Путь каждой черепахи (скорость постоянна по модулю):
\[
S = v t = \frac{2a}{3}.
\]
Встреча происходит в центре исходного треугольника. Траектория каждой черепахи
представляет собой логарифмическую спираль, стягивающуюся к центру.
\end{example}

\begin{example}
\textbf{Собака и лиса.}
Лиса бежит прямолинейно и равномерно со скоростью $u$.
Собака преследует лису со скоростью $v$, вектор которой всегда направлен
на лису. В момент, когда скорости собаки и лисы взаимно перпендикулярны,
расстояние между ними равно $L$.
Найдите ускорение собаки и радиус кривизны её траектории в этот момент.
Обозначим направление от собаки к лисе $\vec{e}_{DF}$.
Скорость собаки направлена вдоль $\vec{e}_{DF}$; скорость лисы
по условию перпендикулярна скорости собаки, то есть перпендикулярна $\vec{e}_{DF}$.
Поперечная компонента скорости лисы (перпендикулярная прямой $DF$):
$u_\perp = u$ (скорость лисы целиком поперечна).
Продольная компонента (вдоль $DF$): $u_\parallel = 0$.
Угловая скорость вращения прямой $DF$:
\[
\omega_{DF} = \frac{u_\perp - 0}{|DF|} = \frac{u}{L}.
\]
Скорость собаки всегда направлена вдоль $DF$, значит, вектор $\vec{v}_{\text{соб}}$
вращается с той же угловой скоростью $\omega_{DF}$.
Ускорение собаки~--- нормальное (скорость постоянна по модулю):
\[
a_{\text{соб}} = v\,\omega_{DF} = \frac{uv}{L}.
\]
Радиус кривизны траектории собаки:
\[
\rho = \frac{v^2}{a_{\text{соб}}} = \frac{vL}{u}.
\]
\end{example}

\noindent\textit{Замечание.} Задачи о погоне второго порядка (преследование по
кривой, когда $v = u$, и пересечение с окружностью) требуют более тонкого анализа
и рекомендуются к самостоятельной проработке.

\end{document}